%%=============================================================================
%% Conclusie
%%=============================================================================

\chapter{Conclusie}%
\label{ch:conclusie}

% TODO: Trek een duidelijke conclusie, in de vorm van een antwoord op de
% onderzoeksvra(a)g(en). Wat was jouw bijdrage aan het onderzoeksdomein en
% hoe biedt dit meerwaarde aan het vakgebied/doelgroep? 
% Reflecteer kritisch over het resultaat. In Engelse teksten wordt deze sectie
% ``Discussion'' genoemd. Had je deze uitkomst verwacht? Zijn er zaken die nog
% niet duidelijk zijn?
% Heeft het onderzoek geleid tot nieuwe vragen die uitnodigen tot verder 
%onderzoek?

In dit onderzoek werd getracht te antwoorden op de vraag hoe Rexx en Python zich verhouden voor automatisatie op z/OS vanuit technisch en praktisch oogpunt. Voor het probleemdomein wordt achterhaald wat de voornaamste redenen zijn voor de modernisatie in mainframes en wat de rol is van automatisatie op mainframes. In automatisatie wordt gekeken naar wat de rol is van Rexx en hoe Rexx de meest gebruikte scripting taal is geworden. Voor het oplossingsdomein worden Python en Rexx met elkaar vergeleken op basis van vooropgestelde requirements, en wordt zo met behulp van meerdere PoC's gekeken of Python nu al een waardige vervanger is van Rexx.

\textbf{Achterhalen wat de voornaamste redenen zijn voor modernisatie in mainframes}

De modernisatie op het mainframe wordt gestimuleerd zowel door een groeiend tekort aan werkkrachten die met mainframes kunnen werken, alsook om te kunnen inspelen op de technologische vooruitgang die wordt gemaakt buiten het mainframe platform. Om de oude technologie op mainframes draaiende te houden zijn nieuwe werkkrachten nodig, en om nieuwe technologie te gebruiken is expertise nodig. Als beide problemen echter op een strategische manier worden aangepakt kan de situatie worden omgedraaid, en kunnen de oplossingen elkaar stimuleren. Als een bedrijf gebruik maakt van nieuwe technologieën, kan het op die manier nieuw talent aantrekken, dat op zijn beurt kan help met het implementeren van nieuwe technologieën op het mainframe, enzovoort. Dit klinkt vanzelfsprekend, maar een probleem dat hiermee nog niet is opgelost, is de nood aan werkkrachten die de huidige systemen kan onderhouden, die (nog) niet worden vernieuwd. Sommige delen van een systeem hebben misschien nog geen moderner alternatief. Een vernieuwing kan ook veel tijd en geld kosten, en in die tijd moeten de systemen wel blijven draaien. Het is dus belangrijk voor bedrijven om een gezonde balans te vinden tussen het blijven moderniseren, om op die manier gebruik te kunnen maken van nieuwe technologieën, en het aantrekken en opleiden van nieuwe talent, zodat een nieuwe werknemer ook een goede achtergrondkennis van de reeds bestaande oudere systemen. In de praktijk zullen beide zaken waarschijnlijk door elkaar lopen.

\textbf{Begrijpen wat automatisatie doet op mainframes en aantonen hoe Rexx daarin de meest gebruikte scripting taal is geworden}

Automatisatie speelt een grote rol op het mainframe. Elke dag worden verschillende batch jobs uitgevoerd die talloze functies hebben, zoals het genereren van rapporten, updaten van databases, en veel meer. Deze batch jobs werken met enorme hoeveelheden data en hebben amper of nooit menselijke interactie nodig. Deze batch jobs worden geschreven in JCL en worden, aan de hand van de functie van de job, meestal aan het begin of het einde van de werkdag uitgevoerd. Grote hoeveelheden worden ook 's nachts uitgevoerd, omdat de druk op het systeem dan meestal het laagst is. Naast de batch job wordt er ook veel gebruik gemaakt van Rexx en Python om scripts te schrijven die allerhande taken uitvoeren. Deze scripts zijn anders dan de JCL scripts, omdat ze meestal worden gebruikt om routinematige taken uit te voeren, die anders manueel door een programmeur zouden worden gedaan. De scripts zelf worden nog steeds uitgevoerd door een programmeur, en kunnen verschillende functies hebben. Voorbeelden van zulke functies zijn het aanmaken van gebruikers, opstarten of afsluiten een systeem, connecteren met API's, enzovoort. De functies van scripts hangen meestal volledig af van de specifieke job van de programmeur.

\textbf{Python en Rexx vergelijken op basis van vooropgestelde requirements, zoals complexiteit en moeilijkheid van onderhoud}

Om Python en Rexx met elkaar te vergelijken is het belangrijk om daar alle achtergrond en nuance bij te nemen. Enkel op die manier kan een volwaardige analyse worden gemaakt van beide talen. Eerst en vooral is het belangrijk te vermelden dat het gebruik Python en Rexx niet de enige twee mogelijkheden zijn om te automatiseren op mainframe. Er zijn andere talen en technologieën, zoals de open source projecten ansible en kubernetes, die een rol kunnen spelen in de automatisatie strategie van bedrijven. Toekomstig onderzoek zou een analyse kunnen maken over het gebruik van die technologieën op het mainframe. Een holistische vergelijking van Python en Rexx kan daarnaast niet zomaar worden gedaan op basis van snelheid en CPU gebruik, of op basis van het aantal mensen die de ene taal over de andere prefereert. De vergelijking van de twee talen gaat veel ruimer dan dat, ook ruimer dan dit onderzoek kan gaan. Om te beginnen moet er een onderverdeling worden gemaakt tussen objectieve en subjectieve aspecten. Objectieve aspecten zoals de performantie, aantal lijnen code, integratie met systemen. Subjectieve aspecten zoals de leer- en leesbaarheid, of algemene gebruiksvriendelijkheid. Een probleem dat is ondervonden in dit onderzoek is dat voor veel van de objectieve aspecten een systeem en onderzoeker nodig zijn die werken zoals in een productie-omgeving. Voor het meten van de performantie (CPU-gebruik, snelheid,...) is een volwaardig systeem nodig. Op het systeem dat werd gebruikt in dit onderzoek zouden de resultaten van een performantie test tot verkeerde conclusies kunnen geleid hebben. Daarnaast is het niet gemakkelijk om subjectieve aspecten zoals leer- en leesbaarheid te meten, aangezien dit meestal meer afhangt van de achtergrond van de programmeur dan van de taal zelf. Dit onderzoek heeft daarom getracht gebruik te maken van enkele specifieke criteria die wel correct konden gemeten worden en die zo min mogelijk afhangen van de ervaring van de programmeur. Die criteria waren nauwer verbonden aan de use cases dan aan de taal in zijn geheel. 

\textbf{Via een vergelijkende studie en PoC's achterhalen of Python nu al een waardige vervanger kan zijn van Rexx, en zo ja, in welke use cases.}